\section{Background and Motivation}
\label{sec:bg}

\subsection{Step Functions and the Amazon States Language}

A Step Functions workflow is a state machine: a \texttt{StartAt} entry and a map of
named \emph{states}. A \texttt{Task} state invokes a Lambda function or a service
integration; \texttt{Choice} branches; \texttt{Parallel} and \texttt{Map} fan out;
\texttt{Wait}, \texttt{Pass}, \texttt{Succeed}, and \texttt{Fail} complete the set.
Data flows through a JSON document that each state reshapes with JSONPath fields
(\texttt{Parameters}, \texttt{InputPath}, \texttt{ResultPath}, \texttt{OutputPath}),
where a key suffixed \texttt{.\$} denotes a path reference. States carry retry
(\texttt{Retry}) and error-handling (\texttt{Catch}) policies. Listing~\ref{lst:asl}
shows a representative task drawn verbatim from our corpus.

\begin{lstlisting}[style=json,caption={A real \asl{} task state (excerpt, \texttt{ship-order}).},label={lst:asl}]
"Get Customer Status": {
  "Type": "Task", "Resource": "arn:aws:states:::lambda:invoke",
  "Parameters": { "FunctionName": "${LambdaGetCustomerStatus}",
                  "Payload.$": "$.detail.customer" },
  "Retry": [ { "ErrorEquals": ["Lambda.ServiceException"], "MaxAttempts": 6 } ],
  "Next": "Do Fraud Check"
}
\end{lstlisting}

The platform validates such a definition's JSON schema and that every transition
target exists, then executes it. It does \emph{not} reason about whether
\texttt{\$.detail.customer} is present, whether retrying this task is safe, or how
the workflow recovers if a later step fails.

\subsection{Four classes of latent defect}
\label{sec:bg:defects}

Our running example is the canonical order-processing workflow
\texttt{CreateOrder} $\to$ \texttt{ReserveInventory} $\to$ \texttt{ChargeCard} $\to$
\texttt{ShipOrder}, which exhibits all four defect classes when edited carelessly. \emph{(1) Contract:}
\texttt{ReserveInventory} produces \texttt{\{orderId, amount\}} but \texttt{ChargeCard}
consumes \texttt{total}---a run-time failure on a missing field. \emph{(2) Typestate:}
reordering to \texttt{ShipOrder\,$\to$\,ChargeCard} ships before payment (valid JSON,
invalid protocol). \emph{(3) Retry:} non-idempotent \texttt{ChargeCard} with a broad
retry (\code{States.ALL}) may charge twice. \emph{(4) Compensation:}
\texttt{ReserveInventory} commits a reservation that leaks if \texttt{ChargeCard} fails
with no release path. None is caught by the platform; all four are recognised,
under-tooled failure modes in cloud-native
systems~\cite{mohammad2025resilient,kulkarni2025characterizing}. We claim nothing about
their prevalence---curated examples are mostly correct
(Section~\ref{sec:eval:wild})---only that they escape static checking entirely today;
Section~\ref{sec:eval} measures detection on both unmodified and injected defects.

\subsection{The information gap}
\label{sec:bg:gap}

The obstacle to checking these properties statically is not the analysis but the
\emph{input}. Raw \asl{} encodes control flow and JSONPath data plumbing, but it does
not record: the \emph{type} of data a task consumes or produces; the \emph{business
state} a task represents; whether a task is \emph{idempotent}; or whether a task is
\emph{persistent} and which action \emph{compensates} it. A verifier for existing
workflows must therefore either (a) restrict itself to properties decidable from
\asl{} alone, or (b) recover the missing semantics. \tool{} does both: it runs sound
native checks that need nothing beyond \asl, and it recovers the remaining semantics
through a lightweight annotation layer whose defaults are inferred from the strong
naming and resource conventions that pervade real workflows
(Section~\ref{sec:approach}).
